Moderná robotika sa v súčasnosti nezaobíde bez komplexných manipulátorov schopných presne, bezpečne a efektívne
interagovať s okolitým prostredím. 
Tento tímový projekt sa zameriava na návrh a implementáciu vnoreného riadiaceho systému pre model robotického 
manipulátora so šiestimi stupňami voľnosti (6 DOF). 
Hlavnou motiváciou je vytvorenie riadiacej jednotky, ktorá úspešne 
spája pokročilé kinematické výpočty s priamym nízkoúrovňovým 
riadením akčných členov v reálnom čase.

Pri interpretácii zadania sme sa zamerali na jednu z najdôležitejších výziev v oblasti robotiky - 
prepočítanie pohybu z priestoru efektora do kĺbových súradníc robota. 
Zatiaľ čo dopredná kinematika nám umožňuje jednoducho určiť polohu a orientáciu koncového efektora na 
základe známych kĺbových uhlov, pre praktické a intuitívne ovládanie je 
kriticky dôležité riešiť inverznú kinematickú úlohu (IKÚ). Bez nej by robotické rameno nedokázalo autonómne 
dosiahnuť konkrétnu cieľovú pozíciu zadanú operátorom. 
Keďže analytické riešenie pre 6-DOF manipulátor môže byť geometricko-matematicky zložité a často 
nejednoznačné, v tomto projekte implementujeme numerické riešenie 
IKÚ pomocou Newtonovej metódy, pričom na efektívne maticové operácie využívame optimalizovanú knižnicu Eigen.
Nepoužívame žiadne externé knižnice špecifické pre robotiku ani ROS, ale vlastnú implementáciu.

Pre realizáciu plne vnoreného (embedded) systému bola zvolená hardvérová platforma mikrokontroléra STM32 
v kombinácii s programovacím jazykom C++. 
Vnorené platformy sú pre riadenie robotických modelov ideálne z dôvodu ich vysokej deterministickosti, 
nízkej latencie, energetickej efektívnosti a možnosti 
priamej hardvérovej integrácie periférií. Celý riadiaci softvér je postavený na operačnom systéme reálneho 
času (RTOS). Tento prístup nám umožnil elegantne rozdeliť 
softvérovú architektúru na samostatné, paralelne bežiace procesy - od periodických výpočtov kinematických 
úloh, cez robustnú medziprocesovú komunikáciu pomocou frontov správ, 
až po spracovanie používateľského rozhrania a samotné riadenie servomotorov.

Okrem samotného algoritmu výpočtu kĺbových premenných IKÚ projekt zahŕňa aj integráciu používateľského rozhrania (LCD displej 
a joystick s tlačidlami) na manuálne zadávanie cieľových súradníc a 
pokročilú nadstavbu v podobe aplikácie 3D mapovania priestoru pomocou laserového senzora vzdialenosti VL53L1X 
s logovaním dát na SD kartu. 
Naším vnútorným cieľom bolo vytvoriť ucelený mechatronický systém, ktorý dokáže spoľahlivo 
fungovať a reagovať aj napriek fyzickým limitom a 
mechanickým vôľam použitého hardvéru.

Pre nás, ako študentov inžinierskeho študijného programu Robotika a kybernetika, predstavuje práca na tomto 
projekte neoceniteľnú skúsenosť. 
Ponúka nám príležitosť prejsť si celým vývojovým cyklom od analýzy geometrie fyzického modelu, cez prenos 
matematických pseudokódov do produkčného C++ kódu, 
až po optimalizáciu a ošetrenie numerickej stability algoritmov v singularitách pomocou regularizácie. 
Tento projekt je vysoko relevantný, 
pretože verne reflektuje reálne priemyselné výzvy v oblasti vývoja embedded softvéru a riadenia robotických 
systémov v reálnom čase.